% Logo of the library: the linear stability regions of the explicit Runge-Kutta
% methods with s = 1..4 stages (from stability_rk<s>.dat, written by
% stability_regions.py; no axes, ticks, labels or legend) with the word "Integrate" in Superclarendon, as a
% violet outline, its "I" filled white at 75% opacity over the curves. The
% word is placed so that the right edge of the stem of its initial "I" lies
% exactly on the (undrawn) Im z axis and the Re z axis crosses the I at
% mid-height. Compile with lualatex (the font is the macOS system font, loaded
% through fontspec).
\documentclass[tikz,border=6pt]{standalone}
\usepackage{fontspec}
\usepackage{pdfrender}
\usepackage{pgfplots}
\pgfplotsset{compat=1.18}

\definecolor{fortranviolet}{RGB}{110,80,146}

\newfontfamily\clarendon{Superclarendon}[UprightFont=* Bold]
\newcommand{\logofont}{\fontsize{120}{120}\clarendon}

% \iedge: distance from the origin of "I" to the right edge of its stem (the
% serifs stick out further), and \imid: height above the baseline of the
% middle of its ink. The stem edge is taken from the outline of the glyph in
% Superclarendon Bold (stem from x = 137 to 340, serifs from 47 to 429, in
% units of 1/1000 em), so it must be updated if the font changes; the height
% is read from the glyph's bounding box. Measured outside the tikzpicture,
% where a real font is current.
\newdimen\iedge
\newdimen\imid
\newcommand{\measureiedge}{\directlua{
  local f = fonts.hashes.identifiers[font.current()]
  local bb = f.descriptions[0x49].boundingbox
  local s = f.parameters.size / f.parameters.units
  local stem = 340 * f.parameters.units / 1000
  tex.setdimen("global", "iedge", math.floor(stem * s + 0.5))
  tex.setdimen("global", "imid", math.floor((bb[2] + bb[4]) / 2 * s + 0.5))
}}

\begin{document}
{\logofont\measureiedge}%
\begin{tikzpicture}
  \begin{axis}[
      x=2.8cm, y=0.625cm,                              % stretched so the I sits inside the RK4 curve
      xmin=-3.2, xmax=1.0, ymin=-3.2, ymax=3.2,        % frame around the four regions
      hide axis,                                       % no axes, only the curves
    ]
    % boundaries |R(z)| = 1, computed by stability_regions.py (one branch per root)
    \addplot[very thick, black!45]           table {stability_rk1.dat};  % Euler: disc of radius 1 around -1
    \addplot[very thick, blue!70!black]      table {stability_rk2.dat};  % two stages, order 2
    \addplot[very thick, orange!80!black]    table {stability_rk3.dat};  % three stages, order 3
    \addplot[very thick, red!70!black]       table {stability_rk4.dat};  % four stages, order 4 (RK4)
    \coordinate (origin) at (axis cs:0,0);                                % where the two axes cross
  \end{axis}

  % ---- the word: the stem of its "I" flush right against the Im z axis, centred on the Re z axis ----
  \begin{scope}[every node/.style={anchor=base west, inner sep=0pt, font=\logofont},
                shift={(origin)}, xshift=-\iedge, yshift=-\imid]
    \node[text opacity=0.75] at (0,0)
         {\textpdfrender{TextRenderingMode=Fill, FillColor=white}{I}};    % white inside of the I, mostly hiding the curves
    \node at (0,0)
         {\textpdfrender{TextRenderingMode=Stroke, LineWidth=1.2pt,      % outline only, as thick as the curves
                         StrokeColor=fortranviolet}
          {Integrate}};                                                   % one node, so the kerning is kept
  \end{scope}
\end{tikzpicture}
\end{document}
